%!TEX program = lualatex
%!TEX encoding = UTF-8 Unicode  
%!TEX root = chapter05.tex  
\documentclass[main.tex]{subfiles} 
\begin{document} 
%----------------------------------------------------------------------------------------
%	CHAPTER 5
%----------------------------------------------------------------------------------------
\cleardoublepage 
\loadgeometry{marginlayout}  
\pagestyle{scrheadings} % Print headers again  
\KOMAoptions{
headwidth=textwithmarginpar,
headsepline=true,
headsepline= 0.5pt,
footwidth=textwithmarginpar,
}   
\onehalfspacing 
\setcounter{chapter}{4}
\chapter{Racine carrée et théorème de Pythagore. Triangles égaux}


\section{La racine carrée}
\begin{definition}
	La racine carrée d'un nombre positif $b\geqslant 0$  est le nombre \emph{positif} noté $\sqrt{b}$ dont le carré vaut $b$. 
	$$
	\left(\sqrt{b}\right)^2=\sqrt{b}\times\sqrt{b} = b $$
	En géométrie $\sqrt{b}$ est \og la longueur du côté d'un carré d'aire  $b$ \fg{}.
\end{definition}
\begin{remark}
	Pour des nombres positifs $b\geqslant 0$, on peut noter $\sqrt{b}=b^{\numprint{0.5}}$. Cette notation est compatible avec les règles d'opérations sur les exposants :
	\begin{align*}
		\big( b^{\numprint{0.5}}\big)^2 & = b^{\numprint{0.5}\times 2} = b^1= b\\
		b^{\numprint{0.5}} \times  b^{\numprint{0.5}} & = b^{\numprint{0.5}+\numprint{0.5}} = b^1 = b
	\end{align*} 
\end{remark}
\begin{theorem} Pour $a$ et $b$ nombres positifs :
	$$\sqrt{a\times b}=\sqrt{a}\times \sqrt{b}$$
\end{theorem}



%----------------------------------------------------------------------------------------
%	EXERCICES
%----------------------------------------------------------------------------------------

\newpage 
\loadgeometry{widelayout}  
\pagestyle{scrheadings} % Print headers again  
\KOMAoptions{
headwidth=textwithmarginpar,
headsepline=true,
headsepline= 0.5pt,
footwidth=textwithmarginpar,
} 
\onehalfspacing
\subsection{Exercices}
estimer des racines carrées.



résoudre des problèmes simples de pythagore


propriétés


\begin{exercise}\hfill\\
	Sans utiliser la calculatrice, calcule (et justifie) :\\
	$\sqrt{25}$\hfill 		
	$\sqrt{7^2}$\hfill 		
	$\sqrt{(3,2)^2}$ \hfill  			
	$\sqrt{(-2,1)^2}$\hfill $\left(\sqrt{4,5}\right)^2$		
	\hfill$\left(\sqrt{-9}\right)^2$
\end{exercise}

\begin{example}[Simplifier] Écrire une expression avec le terme sous la racine le plus petit possible
	
	\parbox[position=t]{0.65\linewidth}{
		$\begin{WithArrows}
			A & = \sqrt{12} \Arrow{identifier le plus grand carré facteur de $12$} \\
			&=\sqrt{4\times 3} \Arrow{utiliser $\sqrt{ab}=\sqrt{a}\sqrt{b}$}\\
			& = \sqrt{4}\sqrt{3}  \\
			& = 2\sqrt{3}  
		\end{WithArrows}$
	}
	\parbox[position=t]{0.3\linewidth}{
		$B = \sqrt{75}$ \hfill $C = 3\sqrt{8}$\\
		\\
		\\
		\\
	}
\end{example}
\begin{exercise}\hfill 
	
	Écrire les expressions suivantes avec le terme sous la racine le plus petit possible :
	\begin{multicols}{5}
		\begin{enumerate}[label={$\Alph*=$}]
			\item $\sqrt{18}$
			\item $\sqrt{20}$
			\item $\sqrt{50}$
			\item $\sqrt{32}$
			\item $\sqrt{200}$
			\item $\sqrt{48}$
			\item $\sqrt{28}$
			\item $\sqrt{99}$
			\item $\sqrt{27}$
			\item $\sqrt{63}$
			\item $\sqrt{45}$
			\item $\sqrt{80}$
			\item $\sqrt{80}$
			\item $\sqrt{150}$
		\end{enumerate}
	\end{multicols}
\end{exercise}
\begin{exercise}\hfill 
	
	Mêmes consignes :
	\begin{multicols}{5}
		\begin{enumerate}[label={$\Alph*=$}]
			\item $5\sqrt{8}$
			\item $2\sqrt{12}$
			\item $3\sqrt{20}$
			\item $4\sqrt{27}$
			\item $2\sqrt{50}$
			\item $5\sqrt{18}$
			\item $6\sqrt{44}$
			\item $5\sqrt{200}$
			\item $3\sqrt{80}$
			\item $10\sqrt{75}$ 
		\end{enumerate}
	\end{multicols}
\end{exercise}
\begin{example} Sans utiliser de calculatrice, comparer $2\sqrt{3}$ et $3\sqrt{2}$ :
	
	\vfill 
\end{example}

\begin{exercise}\hfill \\
	Écrire les nombres suivants sous la forme $\sqrt{k}$ ou $k$ est un entier. Par exemple $2\sqrt{2}=\sqrt{8}$.  
	$$ 3\sqrt{5},\quad 6\sqrt{3}; \quad 8\sqrt{5}; \quad 4\sqrt{6}; \quad 10\sqrt{10}; \quad 7\sqrt{7} $$
\end{exercise}
\begin{exercise}\hfill  
	\begin{enumerate}[label={\alph*)},leftmargin=*]
		\item Quelle est l'aire $\mathscr{A}$ d'un carré de côté $3^7~\si{\centi\meter}$ ?
		\item Combien mesure le côté d'un carré d'aire $3^{10}~\si{\centi\meter^2}$ ? 
		\item Quelle est l'aire d'un triangle rectangle isocèle dont les côtés de l'angle droite mesurent $5^3~\si{\milli\meter}$ ? 
		\item Quel est le périmètre du triangle précédent ?
		\item Que devient l'aire de mon triangle précédent, si je multiplie les longueurs par 3? 
	\end{enumerate}
\end{exercise} 


\begin{exercise} \hfill  \\ 
	\parbox{0.55\linewidth}{
		\begin{enumerate}[label=\alph*),leftmargin=*]
			\item Donne la longueur exacte du côté $[CA]$, puis calcule une valeur approchée au centimètre prés.
			\item Le triangle $ACD$ est-il rectangle?  Justifie.
		\end{enumerate}
	}\hfill 
	\parbox{0.35\linewidth}{
		\begin{tikzpicture}[scale=0.6, line cap=round,line join=round,>=triangle 45,x=1cm,y=1cm]
			\draw [line width=1pt] (3 ,-0.5) coordinate(A) node[below]{$A$} -- node[midway,below]{\scriptsize $7$ cm} (5.14,-0.56);
			\draw [line width=1pt] (5.14,-0.5) coordinate(B) node[below]{$B$}--  node[midway,right]{\scriptsize $14$ cm}(3,3) coordinate(C) node[above]{$C$};
			\draw [line width=1pt] (C)-- (A);
			\draw[color=black] (-4,3) coordinate(D) node[left]  {$D$};
			\draw [line width=1pt] (C)-- node[midway,above]{\scriptsize $35$ cm} (D);
			\draw [line width=1pt] (A)-- node[midway,below]{\scriptsize $37$ cm} (D);
			% angle droit
			\draw (3,0)-|(3.5,-0.5);
		\end{tikzpicture}
	}
\end{exercise}
\begin{exercise}\hfill  \\%12 points
	\parbox{0.65\linewidth}{La figure ci-contre est codée et réalisée à main levée.
		
		Elle représente un quadrilatère $ABCD$ dont les
		diagonales se croisent en un point $O$. On donne: OA $= 3,5$ cm et AB $= 5$ cm.
		
		On s'intéresse à la nature du quadrilatère ABCD qui a été représenté. 
		\begin{enumerate}[label=\arabic*),leftmargin=*]
			\item Peut-on affirmer que ABCD est un rectangle ? Justifier.
			\item Peut-on affirmer que ABCD est un carré ? Justifier.
		\end{enumerate} 
	}\hfill
	\parbox{0.3\linewidth}{
		\begin{tikzpicture}% [scale=0.5]
			\tkzDefPoint(0.2,3.8){A}
			\tkzDefPoint(3.7,3.6){B}
			\tkzDefPoint(3.6,0.5){C}
			\tkzDefPoint(0.5,0.5){D}
			\tkzDrawSegments[decorate, decoration={snake, amplitude=.3mm}](A,B A,C B,C C,D A,D B,D)
			
			\pgfresetboundingbox % removes white space 
			\tkzSetUpPoint[size=3, fill=black]
			\tkzfctset{tan style/.style={->,>=latex, color=red, line width=1.2pt}} 
			%	 		\tkzInit[ymin=-5,ymax=5.5,xmin=-1,xmax=25]
			\tkzDrawPoints(A,B,C,D)  
			\tkzLabelPoints[above left](A)
			\tkzLabelPoints[above right](B) 
			\tkzLabelPoints[below right](C)
			\tkzLabelPoints[below left](D)
			\tkzMarkSegments[pos=.25,mark=s||](A,C B,D);
			\tkzMarkSegments[pos=.75,mark=s||](A,C B,D);
			\tkzDefPoint(2,2.1){O}
			\tkzLabelPoints[left](O)
			%	 		\tkzLabelSegment[midway, sloped, below](A,T){ \scriptsize 2,8 \si{\centi\meter}};  
		\end{tikzpicture} 
	} 
	
\end{exercise}
\begin{exercise}[\cancel{\faCalculator}]\hfill 


\begin{tikzpicture}[scale=0.85]
%	\tkzInit[xmin=-1.0,ymin=-1,xmax=6,ymax=3.5]
%	\tkzClip
	\tkzSetUpPoint[size=3, fill=black]
	\tkzfctset{tan style/.style={->,>=latex, color=red, line width=1.2pt}} 
	\tkzDefPoint(0,0){E} \tkzDefPoint(4.66,0){D}
	\tkzDefMidPoint(E,D) \tkzGetPoint{p}
	\tkzInterCC[R](p,2.33cm)(E,4cm)\tkzGetPoints{m}{C}
	\tkzDrawPolygon(E,C,D)
	\tkzLabelSegment[below](E,D){\scriptsize$8$}
	\tkzLabelSegment[above left](E,C){\scriptsize$x$}
	\tkzLabelSegment[above right](D,C){\scriptsize$2\sqrt{2}$}
	%	\tkzDrawPoints(E,C,D)
	%	\tkzLabelPoints[below left](E)
	%	\tkzLabelPoints[below right](D)
	%	\tkzLabelPoints[above ](C) 
	%		\tkzLabelAngle[pos=3.5,circle](N,P,M){?} 
	\tkzMarkRightAngle[ size=0.2](E,C,D)
\end{tikzpicture} 
\begin{tikzpicture}[scale=0.80 ]
\tkzInit[xmin=-2.5 ,ymin=-0.75,xmax=4.25,ymax=3.5]
\tkzClip
\tkzSetUpPoint[size=7, fill=black]
\tkzfctset{tan style/.style={->,>=latex, color=red, line width=1.2pt}} 
\tkzDefPoint(0,0){A} 
%		\tkzDefPoint(0,2.33){B} 
\tkzDefShiftPointCoord[A](145:2.33){B}
%		 \tkzDefPoint(4,0){C}
\tkzDefShiftPointCoord[A](55:4){C}
\tkzDrawPolygon(B,A,C)
%		\tkzLabelPoints[below left](A)
%		\tkzLabelPoints[below right](C)
%		\tkzLabelPoints[above left](B)
\tkzLabelSegment[  below, sloped](B,A){\scriptsize$7$}
\tkzLabelSegment[below, sloped](A,C){\scriptsize$\sqrt{11}$} 
\tkzLabelSegment[above , sloped](B,C){\scriptsize$y$} 
\tkzMarkRightAngle[  size=0.2](B,A,C)
\end{tikzpicture}  
\begin{tikzpicture}[scale=0.80]
	\tkzInit[xmin=-0.5 ,ymin=-0.75,xmax=4.75,ymax=3.5]
	\tkzClip
	\tkzSetUpPoint[size=7, fill=black]
	\tkzfctset{tan style/.style={->,>=latex, color=red, line width=1.2pt}} 
	\tkzDefPoint(0,0){A} \tkzDefPoint(0,2.33){B} \tkzDefPoint(4,0){C}
	\tkzDrawPolygon(B,A,C)
	%		\tkzLabelPoints[below left](A)
	%		\tkzLabelPoints[below right](C)
	%		\tkzLabelPoints[above left](B)
	\tkzLabelSegment[left](B,A){\scriptsize$z$} 
	\tkzLabelSegment[below](A,C){\scriptsize$\sqrt{10}$}
	\tkzLabelSegment[above right](B,C){\scriptsize$10$} 
	\tkzMarkRightAngle[  size=0.2](B,A,C)
\end{tikzpicture}  
\begin{tikzpicture}[scale=0.80 ]
			\tkzInit[xmin=-1.5 ,ymin=-0.75,xmax=4.25,ymax=3.5]
			\tkzClip
			\tkzSetUpPoint[size=7, fill=black]
			\tkzfctset{tan style/.style={->,>=latex, color=red, line width=1.2pt}} 
			\tkzDefPoint(0,0){A} 
	%		\tkzDefPoint(0,2.33){B} 
			\tkzDefShiftPointCoord[A](101:2.1){B}
	%		 \tkzDefPoint(4,0){C}
			 \tkzDefShiftPointCoord[A](10:4){C}
			\tkzDrawPolygon(B,A,C)
					\tkzLabelPoints[below left](A)
					\tkzLabelPoints[below right](C)
					\tkzLabelPoints[above left](B)
			\tkzLabelSegment[  below, sloped](B,A){\scriptsize$4\sqrt{5}$}
			\tkzLabelSegment[below, sloped](A,C){\scriptsize$2\sqrt{5}$} 
			\tkzLabelSegment[above , sloped](B,C){\scriptsize$10$} 
			\tkzMarkRightAngle[  size=0.2](B,A,C)
		\end{tikzpicture} 
	\begin{enumerate}[label=\alph*),leftmargin=*]
		\item Calculer la longueur manquante et donner le résultat sous la forme $a\sqrt{k}$, avec $k$ entier le plus petit possible.
		\item Justifier soigneusement que le triangle $ABC$ est rectangle.
	\end{enumerate}
\end{exercise}
\begin{exercise}[\cancel{\faCalculator}]\hfill 
 
\noindent\parbox{0.55\linewidth}{
Soit le pavé droit $ABCDEFGH$ ci-contre :
	\begin{enumerate}
		\item Montrer que  $AC=5$.
		\item Calculer $AG$.
	\end{enumerate}
}\hfill
\parbox{0.4\linewidth}{\centering 
		\begin{tikzpicture}[scale=0.55] 
			\def \largeur{6};
			\def \hauteur{4}; 
			\def \profondeur{4}; 
			\def \angle{30}; 
			\pgfmathparse{\profondeur*cos(\angle)/2}\let\x\pgfmathresult
			\pgfmathparse{\profondeur*sin(\angle)/2}\let\y\pgfmathresult
			\pgfmathparse{\largeur+\x}\let\xx\pgfmathresult
			\pgfmathparse{\hauteur+\y}\let\yy\pgfmathresult
			
			\draw[  color=black,opacity=1] (0,0) -- (\largeur,0) -- (\largeur,\hauteur) -- (0,\hauteur) -- cycle;
			\draw[dashed,black] (0,0) -- (\x,\y) -- (\xx,\y); 
			\draw[dashed,black] (\x,\y) -- (\x,\yy);
			\draw[ black,opacity=1] (0,\hauteur) -- (\x,\yy) -- (\xx,\yy) -- (\largeur,\hauteur) -- cycle;
			\draw[black,opacity=1] (\largeur,0) -- (\largeur,\hauteur) -- (\xx,\yy) -- (\xx,\y) -- cycle;
			
			\draw[<->,>=stealth',black, dashed] (0,0) --  (\xx,\yy); 
			
			\draw[<->,>=stealth',black, dashed] (0,0) --  (\xx,\y); 
			
			\node[below left,black] (A) at (0,0) {\scriptsize A};
			\node[below right, black] (B) at (\largeur,0) {\scriptsize B};
			\node[below right,black] (C) at (\xx,\y) {\scriptsize C};
			\node[below,black] (D) at (\x,\y) {\scriptsize D};
			\node[left,black] (E) at (0,\hauteur) {\scriptsize E};
			\node[below right,black] (F) at (\largeur,\hauteur) {\scriptsize F};
			\node[right,black] (G) at (\xx,\yy) {\scriptsize G};
			\node[above,black] (H) at (\x,\yy) {\scriptsize H};
			
			\draw[<->,>=stealth'] (A) -- (B) node[midway,below] {$3\sqrt{2}$};
			\draw[<->,>=stealth'] (B) -- (C) node[midway,below right] {$\sqrt{7}$};
			\draw[<->,>=stealth'] (C) -- (G) node[midway,right] {$2\sqrt{5}$};
			
			%\begin{scope}[xshift =8 cm , yshift =+1 cm, scale=1.0]
			%\patronpave[pos =1 , codages ,  a =1 , b =2 , c =3 ]
			%\end{scope}
		\end{tikzpicture}
	} 
	\end{exercise}
\begin{exercise} \hfill
	
\noindent Pour le triangle $ABC$ ci-dessous, les longueurs des côtés sont exprimées en fonction d'un nombre $x$ :\\
	\begin{minipage}[T]{0.4\columnwidth} \centering
		$
		\begin{cases}
			AB=3x+6 \\
			BC=4x+8 \\
			AC= 5x+10 \end{cases}$
	\end{minipage}
	\hfill 
	\begin{minipage}[T]{0.50\columnwidth}  
		\begin{tikzpicture}[scale=0.3]
			\tkzInit[xmin=-3.5,ymin=-2.0,xmax=17,ymax=6.0]
			\tkzClip
			\tkzDefPoint(0,0){B} \tkzDefPoint(15,0){C}
			\tkzDrawTriangle[two angles = 90 and 15](B,C)\tkzGetPoint{A}
			\tkzDrawPoints[color=black](A,B,C)
			\tkzLabelPoints[right](C)
			\tkzLabelPoints[left](A,B)
		\end{tikzpicture} 
	\end{minipage}
\begin{enumerate}[label={\alph*)},leftmargin=*]
	\item Faire la figure lorsque $x=0$. Le triangle obtenu est-il rectangle ? \textbf{Justifie}.
	\item Montrer que pour tout nombre $x$ positif, le triangle $ABC$ est rectangle en B.
\end{enumerate}
\end{exercise}
%----------------------------------------------------------------------------------------
%	 
%----------------------------------------------------------------------------------------
\newpage
\loadgeometry{marginlayout}  
\pagestyle{scrheadings} % Print headers again  
\KOMAoptions{
	headwidth=textwithmarginpar,
	headsepline=true,
	headsepline= 0.5pt,
	footwidth=textwithmarginpar,
}   
\onehalfspacing 
\section{Égalité des triangles}


\begin{definition}[triangles égaux]
	Deux triangles sont égaux si leurs trois côtés et leur trois angles sont égaux deux à deux. 
\end{definition} 
Les triangles $ABC$ et $IJK$ de la figure \ref{chapitre05:trianglesegaux} sont égaux. On a les 6 égalités suivantes :
\begin{figure}[h]\centering
	\begin{tikzpicture}[scale=0.90]
		\tkzDefPoint(0,0){A}
		\tkzLabelPoint[above left](A){$A$}
		\tkzDefShiftPoint[A](-135:2.5){B}  
		\tkzDefTriangle[two angles = 80 and 60](A,B)
		\tkzGetPoint{C}
		\tkzLabelPoint[left](B){$B$}
		\tkzLabelPoint[right](C){$C$}
		\tkzDrawPolygon(A,B,C)
		\tkzDrawPoints(A,B,C)
		\tkzDefShiftPoint[A](4,+3.0){J}  
		\tkzDefShiftPoint[J](-70:2.5){I}
		\tkzDefTriangle[two angles = -80  and -60](I,J)
		\tkzGetPoint{K}
		\tkzDrawPolygon(I,J,K)
		\tkzDrawPoints(I,J,K) 
		\tkzLabelPoint[right](I){$I$}
		\tkzLabelPoint[left](J){$J$}
		\tkzLabelPoint[right](K){$K$}
		\tkzMarkAngle[mark=|,mkcolor=black, thick, size=0.6, fill=red!20, opacity=0.6](B,A,C)
		\tkzMarkAngle[mark=||,mkcolor=black, thick,size=0.6, fill=blue!20, opacity=0.6](J,K,I)
		\tkzMarkAngle[mark=||,mkcolor=black, thick,size=0.6, fill=blue!20, opacity=0.6](A,C,B)
		\tkzMarkAngle[mark=|,mkcolor=black, thick,size=0.6, fill=red!20, opacity=0.6](K,I,J)
		\tkzMarkAngle[mark=s,mkcolor=black, thick,size=0.6, fill=green!20, opacity=0.6](C,B,A)
		\tkzMarkAngle[mark=s,mkcolor=black, thick,size=0.6, fill=green!20, opacity=0.6](I,J,K) 
		\tkzDrawSegment[line width=2pt, color=red](B,C)
		\tkzDrawSegment[line width=2pt, color=red](J,K)
		\tkzDefMidPoint(B,C)\tkzGetPoint{D}
		\tkzDefMidPoint(J,K)\tkzGetPoint{L}
		\draw[<->, >=latex] (A) .. controls +(right:1cm) and +(left:2cm) .. (I) node[sloped,above,pos=0.5]{\scriptsize Angles homologues} ;
		\draw[<->, >=latex] (D) .. controls +(right:8cm) and +(down:4cm) .. (L) node[ below right,pos=0.6]{\scriptsize Côtés homologues} ;
	\end{tikzpicture} 
	\caption{$ABC$ et $IJK$ sont égaux.
%		\\
%		${\color{red} \widehat{A}= \widehat{I} }$\\
%		$ {\color{ForestGreen} \widehat{B}= \widehat{J} } $\\
%		$ {\color{Blue} \widehat{C}= \widehat{K} }$\\
%		$ {\color{red} BC=JK }$\\
%		$ {\color{ForestGreen}  AC=IK } $\\
%		$ {\color{Blue} AB=IJ}$ 
} \label{chapitre05:trianglesegaux}
\end{figure}   
\setlength{\abovedisplayskip}{0pt}
\setlength{\belowdisplayskip}{0pt} 
\begin{align*}
	& {\color{red} \widehat{A}= \widehat{I} }
	&&  {\color{ForestGreen} \widehat{B}= \widehat{J} } 
	&& {\color{Blue} \widehat{C}= \widehat{K} } \\
	& {\color{red} BC=JK }
	&&  {\color{ForestGreen}  AC=IK } 
	&& {\color{Blue} AB=IJ} 
\end{align*} 
Il n'est pas nécessaire de vérifier les six égalités pour prouver que les triangles sont égaux et que les 6 égalités sont vraies.  
Il suffit de vérifier que l'on est dans l'une des trois situations suivantes.
 
\begin{postulat}[Critère CCC] 
	Si deux triangles ont leurs trois côtés respectivement égaux, alors ils sont égaux.
	\begin{center}
		\begin{tikzpicture}
			\tkzDefPoint(0,0){A}
			%\tkzLabelPoint[above left](A){$A$}
			\tkzDefShiftPoint[A](-135:2.5){B}  
			\tkzDefTriangle[two angles = 80 and 60](A,B)
			\tkzGetPoint{C}
			%\tkzLabelPoint[left](B){$B$}
			%\tkzLabelPoint[right](C){$C$}
			\tkzDrawPolygon(A,B,C)
			\tkzDrawPoints(A,B,C)
			\tkzDefShiftPoint[A](7.0,-2.0){J}  
			\tkzDefShiftPoint[J](110:2.5){I}
			\tkzDefTriangle[two angles = -80  and -60](I,J)
			\tkzGetPoint{K}
			\tkzDrawPolygon(I,J,K)
			\tkzDrawPoints(I,J,K) 
			%\tkzLabelPoint[right](I){$I$}
			%\tkzLabelPoint[below left](J){$J$}
			%\tkzLabelPoint[below](K){$K$}
			\tkzMarkSegments[mark=s|, color=black, thick](B,C K,J)
			\tkzMarkSegments[mark=s||, color=black, thick](B,A I,J)
			\tkzMarkSegments[mark=x,color=black, thick](A,C K,I)
			%\tkzMarkAngle[mark=|,mkcolor=black, thick, size=0.6, fill=red!20, opacity=0.6](B,A,C)
			%\tkzMarkAngle[mark=||,mkcolor=black, thick,size=0.6, fill=blue!20, opacity=0.6](J,K,I)
			%\tkzMarkAngle[mark=||,mkcolor=black, thick,size=0.6, fill=blue!20, opacity=0.6](A,C,B)
			%\tkzMarkAngle[mark=|,mkcolor=black, thick,size=0.6, fill=red!20, opacity=0.6](K,I,J)
			%\tkzMarkAngle[mark=s,mkcolor=black, thick,size=0.6, fill=green!20, opacity=0.6](C,B,A)
			%\tkzMarkAngle[mark=s,mkcolor=black, thick,size=0.6, fill=green!20, opacity=0.6](I,J,K) 
			\tkzDrawSegment[line width=1pt ](B,C)
			\tkzDrawSegment[line width=1pt ](J,K)
			\tkzDefMidPoint(B,C)\tkzGetPoint{D}
			\tkzDefMidPoint(J,K)\tkzGetPoint{L}
		\end{tikzpicture}
	\end{center}
\end{postulat} 
\begin{remark}
	Cas RHC (Rectangle-Hypoténuse-Côté). Deux triangles rectangles, qui ont même longueur d'hypoténuse et une même longueur d'un côté de l'angle droit sont égaux.
\end{remark}

\newpage 
\begin{postulat}[Critère CAC]   Si deux triangles ont un angle égal compris entre deux côtés respectivement égaux, alors ils sont égaux.
	\begin{center}
		\begin{tikzpicture}
			\tkzDefPoint(0,0){A}
			%\tkzLabelPoint[above left](A){$A$}
			\tkzDefShiftPoint[A](-125:2.5){B}  
			\tkzDefTriangle[two angles = 80 and 60](A,B)
			\tkzGetPoint{C}
			%\tkzLabelPoint[left](B){$B$}
			%\tkzLabelPoint[right](C){$C$}
			\tkzDrawPolygon(A,B,C)
			\tkzDrawPoints(A,B,C)
			\tkzDefShiftPoint[A](7.0,-2.0){J}  
			\tkzDefShiftPoint[J](110:2.5){I}
			\tkzDefTriangle[two angles = -80  and -60](I,J)
			\tkzGetPoint{K}
			\tkzDrawPolygon(I,J,K)
			\tkzDrawPoints(I,J,K) 
			%\tkzLabelPoint[right](I){$I$}
			%\tkzLabelPoint[below left](J){$J$}
			%\tkzLabelPoint[below](K){$K$}
			%\tkzMarkSegments[mark=x,mkcolor=black, thick](B,C K,J)
			\tkzMarkSegments[mark=s|,color=black, thick](B,A I,J)
			\tkzMarkSegments[mark=s||,color=black, thick](A,C K,I)
			\tkzMarkAngle[mark=s,color=black, thick, size=0.6, fill=red!20, opacity=0.6](B,A,C)
			%\tkzMarkAngle[mark=||,mkcolor=black, thick,size=0.6, fill=blue!20, opacity=0.6](J,K,I)
			%\tkzMarkAngle[mark=||,mkcolor=black, thick,size=0.6, fill=blue!20, opacity=0.6](A,C,B)
			\tkzMarkAngle[mark=s,mkcolor=black, thick,size=0.6, fill=red!20, opacity=0.6](K,I,J)
			%\tkzMarkAngle[mark=s,mkcolor=black, thick,size=0.6, fill=green!20, opacity=0.6](C,B,A)
			%\tkzMarkAngle[mark=s,mkcolor=black, thick,size=0.6, fill=green!20, opacity=0.6](I,J,K) 
			\tkzDrawSegment[line width=1pt ](B,C)
			\tkzDrawSegment[line width=1pt ](J,K)
			\tkzDefMidPoint(B,C)\tkzGetPoint{D}
			\tkzDefMidPoint(J,K)\tkzGetPoint{L}
		\end{tikzpicture}
	\end{center}
	
\end{postulat} 

\begin{postulat}[Critère ACA] Si deux triangles ont un côté égal adjacent à deux angles respectivement égaux, alors ils sont égaux.
	\begin{center}
		\begin{tikzpicture}
			\tkzDefPoint(0,0){A}
			%\tkzLabelPoint[above left](A){$A$}
			\tkzDefShiftPoint[A](-135:2.5){B}  
			\tkzDefTriangle[two angles = 80 and 60](A,B)
			\tkzGetPoint{C}
			%\tkzLabelPoint[left](B){$B$}
			%\tkzLabelPoint[right](C){$C$}
			\tkzDrawPolygon(A,B,C)
			\tkzDrawPoints(A,B,C)
			\tkzDefShiftPoint[A](7.0,-2.0){J}  
			\tkzDefShiftPoint[J](110:2.5){I}
			\tkzDefTriangle[two angles = -80  and -60](I,J)
			\tkzGetPoint{K}
			\tkzDrawPolygon(I,J,K)
			\tkzDrawPoints(I,J,K) 
			%\tkzLabelPoint[right](I){$I$}
			%\tkzLabelPoint[below left](J){$J$}
			%\tkzLabelPoint[below](K){$K$}
			\tkzMarkSegments[mark=s||,color=black, thick](B,C K,J)
			%\tkzMarkSegments[mark=x,mkcolor=black, thick](B,A I,J)
			%\tkzMarkSegments[mark=s|,mkcolor=black, thick](A,C K,I)
			%\tkzMarkAngle[mark=||,mkcolor=black, thick, size=0.6, fill=red!20, opacity=0.6](B,A,C)
			\tkzMarkAngle[mark=|,mkcolor=black, thick,size=0.6, fill=blue!20, opacity=0.6](J,K,I)
			\tkzMarkAngle[mark=|,mkcolor=black, thick,size=0.6, fill=blue!20, opacity=0.6](A,C,B)
			%\tkzMarkAngle[mark=||,mkcolor=black, thick,size=0.6, fill=red!20, opacity=0.6](K,I,J)
			\tkzMarkAngle[mark=s,mkcolor=black, thick,size=0.6, fill=green!20, opacity=0.6](C,B,A)
			\tkzMarkAngle[mark=s,mkcolor=black, thick,size=0.6, fill=green!20, opacity=0.6](I,J,K) 
			\tkzDrawSegment[line width=1pt ](B,C)
			\tkzDrawSegment[line width=1pt ](J,K)
			\tkzDefMidPoint(B,C)\tkzGetPoint{D}
			\tkzDefMidPoint(J,K)\tkzGetPoint{L}
		\end{tikzpicture}
	\end{center}
\end{postulat}
 

\begin{remark} Il n'y a pas de critère ACC
	% Tracer deux triangles non égaux, ayant un angle de \ang{40}, adjacent à un côté de \numprint{4.5}~\si{\centi\meter} et opposé à un côté de \numprint{3.5}~\si{\centi\meter} de longueur.
	\begin{figure}[h]\centering
		\begin{tikzpicture}[scale=1.0]
			\tkzDefPoint(0,0){A}
			\tkzDefShiftPoint[A](40:5){B}
			\tkzDefMidPoint(A,B)\tkzGetPoint{I}
			\tkzDefShiftPoint[I](-0.2,0.2){I'}
			\node[inner sep=0pt, opacity=1,   rotate=-15] (carte) at (I') {
				\includegraphics[ width=0.25\linewidth]{fishing-pole2}
			}; 
			\tkzDefShiftPoint[A](0:6){C}
			%\tkzDrawPoint(B)
			\tkzDrawLine[add= 0.2 and 0.25, dashed](A,C)
			\tkzDrawLine[add= 0 and 0.15, dashed](A,B)
			\tkzLabelAngle[pos=1.5](C,A,B){\large\ang{40}} 
			\tkzMarkAngle[size=1.0, opacity=0.5,mark=none](C,A,B)
		\end{tikzpicture}
		\caption{}
	\end{figure} 
\end{remark}  
\begin{remark} Avoir 2 angles homologues égaux n'est pas suffisant pour dire que les triangles sont égaux. 
	\end{remark}
%----------------------------------------------------------------------------------------
%	EXERCICES
%----------------------------------------------------------------------------------------

\newpage 
\loadgeometry{widelayout}  
\pagestyle{scrheadings} % Print headers again  
\KOMAoptions{
	headwidth=textwithmarginpar,
	headsepline=true,
	headsepline= 0.5pt,
	footwidth=textwithmarginpar,
} 
\onehalfspacing
\subsection{Exercices}


\begin{exercise}[Égalité des triangles]\hfill 
	
	Si possible, démontrer que les paires de triangles sont égaux. Les figures ne sont pas à l'échelle.
\end{exercise} 
\noindent
%\fbox{
\parbox[position=t]{0.48\textwidth}{
	%\begin{figure}[H]\centering
	\begin{tikzpicture}[scale=1.0]
		\tkzText(0.5,2){$(A)$}
		\tkzSetUpPoint[size=3, fill=black]
		\tkzfctset{tan style/.style={->,>=latex, color=red, line width=1.2pt}}   
		\tkzSetUpLine[line width= 1.2pt, add=.5 and .5] 
		\tkzDefPoint(0,0){A}
		\tkzDefPoint(3,0){B}
		\tkzDefShiftPoint[A](-26:3.3){B}
		\tkzDefShiftPoint[A](32:4.5){C} 
		\tkzDrawPolygon[line width=1.2pt](A,B,C)
		%		\tkzLabelPoints[font=\scriptsize](A,B,C)
		\tkzMarkAngle[mark=none,-, >=latex, size =0.5](C,B,A)
		\tkzMarkAngle[mark=none,-, >=latex, size =0.5](A,C,B)
		\tkzLabelAngle[font=\scriptsize,pos=.8](C,B,A){$\ang{70}$}
		\tkzLabelAngle[font=\scriptsize,pos=.8](A,C,B){$\ang{55}$}
		\tkzLabelSegment[above left, font=\scriptsize](A,C){\SI{10}{\centi\meter}}
		%		\tkzDefPointWith[orthogonal,K= -0.65](C,A)
		%		\tkzGetPoint{D}
		%		\tkzDefLine[parallel=through D](A,B) \tkzGetPoint{E}
		%		\tkzDrawLine(D,E)
		%		\tkzDrawLine(A,B)
		%		\tkzLabelLine[pos=1.5,  right](A,B){$(\Delta_1)$}
		%		\tkzLabelLine[pos=1.5,  right](D,E){$(\Delta_2)$}
		%		\tkzMarkRightAngle(D,C,A)
		%		
		%		\tkzDrawSegments[line width= 1.2pt](A,C C,D)
		%		\tkzMarkAngle[mark=none,-, >=latex, size =0.7](B,A,C)
		%		\tkzLabelAngles[pos=1.3,  ](B,A,C){$a$}
		%		\tkzMarkAngle[mark=none,-, >=latex, size =0.7](C,D,E)
		%		\tkzLabelAngles[pos=1.3, ](C,D,E){$b$} 
		\begin{scope}[xshift=4cm] 
			\tkzDefPoint(0,0){A}
			\tkzDefPoint(3,0){B}
			\tkzDefShiftPoint[A](-5:4.1){B}
			\tkzDefShiftPoint[A](48:3.3){C} 
			\tkzDrawPolygon[line width=1.2pt](A,B,C)
			%	\tkzLabelPoints[font=\scriptsize](A,B,C)
			\tkzMarkAngle[mark=none,-, >=latex, size =0.5](B,A,C)
			\tkzMarkAngle[mark=none,-, >=latex, size =0.5](A,C,B)
			\tkzLabelAngle[font=\scriptsize,pos=.8](B,A,C){$\ang{55}$}
			\tkzLabelAngle[font=\scriptsize,pos=.8](A,C,B){$\ang{70}$} 
			\tkzLabelSegment[below, font=\scriptsize](A,B){\SI{10}{\centi\meter}}
		\end{scope}
	\end{tikzpicture}   
}
%}
\hfill 
%\fbox{
\parbox[position=t]{0.45\textwidth}{
	%\begin{figure}[H]\centering
	\begin{tikzpicture}[scale=1.0]
		\tkzText(0.5,2.8){$(B)$}
		\tkzSetUpPoint[size=3, fill=black]
		\tkzfctset{tan style/.style={->,>=latex, color=red, line width=1.2pt}}   
		\tkzSetUpLine[line width= 1.2pt, add=.5 and .5] 
		\tkzDefPoint(1,0){A} 
		\tkzDefShiftPoint[A](-5:3.3){B}
		\tkzDefTriangle[two angles=90 and 35](A,B)\tkzGetPoint{C}  
		\tkzDrawPolygon[line width=1.2pt](A,B,C)
		%		\tkzLabelPoints[font=\scriptsize](A,B,C)
		\tkzMarkRightAngle[](B,A,C)
		%\tkzMarkAngle[mark=none,-, >=latex, size =0.5](C,B,A)
		%\tkzMarkAngle[mark=none,-, >=latex, size =0.5](A,C,B)
		%\tkzLabelAngle[font=\scriptsize,pos=.8](C,B,A){$\ang{70}$}
		%\tkzLabelAngle[font=\scriptsize,pos=.8](A,C,B){$\ang{55}$}
		%\tkzLabelSegment[above, font=\scriptsize](A,C){$10$} 
		\tkzLabelSegment[below, font=\scriptsize](A,B){\SI{7}{\centi\meter}}
		\tkzLabelSegment[above, font=\scriptsize](C,B){\SI{9}{\centi\meter}}
		\begin{scope}[xshift=5cm] 
			\tkzDefPoint(0,0){A}
			\tkzDefShiftPoint[A](+15:2.5){B} 
			\tkzDefTriangle[two angles=90 and 35](A,B)\tkzGetPoint{C}  
			
			\tkzDrawPolygon[line width=1.2pt](A,B,C)
			%	\tkzLabelPoints[font=\scriptsize](A,B,C)
			\tkzMarkRightAngle[](B,A,C)
			%\tkzMarkAngle[mark=none,-, >=latex, size =0.5](C,B,A)
			%\tkzMarkAngle[mark=none,-, >=latex, size =0.5](A,C,B)
			%\tkzLabelAngle[font=\scriptsize,pos=.8](C,B,A){$\ang{70}$}
			%\tkzLabelAngle[font=\scriptsize,pos=.8](A,C,B){$\ang{55}$}
			%\tkzLabelSegment[above, font=\scriptsize](A,C){$10$} 
			\tkzLabelSegment[left, font=\scriptsize](A,C){\SI{7}{\centi\meter}}
			\tkzLabelSegment[above right, font=\scriptsize](C,B){\SI{9}{\centi\meter}}
		\end{scope}
	\end{tikzpicture}   
}
%}

\noindent
%\fbox{
\parbox[position=t]{0.45\textwidth}{
	%\begin{figure}[H]\centering
	\begin{tikzpicture}[scale=1.0]
		\tkzText(0.5,2.8){$(C)$}
		\tkzSetUpPoint[size=3, fill=black]
		\tkzfctset{tan style/.style={->,>=latex, color=red, line width=1.2pt}}   
		\tkzSetUpLine[line width= 1.2pt, add=.5 and .5] 
		\tkzDefPoint(0.5,0){A} 
		\tkzDefShiftPoint[A](-5:2.8){B}
		\tkzDefTriangle[two angles=55 and 80](A,B)\tkzGetPoint{C}  
		\tkzDrawPolygon[line width=1.2pt](A,B,C)
		%		\tkzLabelPoints[font=\scriptsize](A,B,C)
		%		\tkzMarkRightAngle[](B,A,C)
		\tkzMarkAngle[mark=none,-, >=latex, size =0.5](B,A,C)
		\tkzMarkAngle[mark=none,-, >=latex, size =0.5](C,B,A)
		\tkzMarkAngle[mark=none,-, >=latex, size =0.5](A,C,B)
		\tkzLabelAngle[font=\scriptsize,pos=.8](B,A,C){$\ang{55}$}
		\tkzLabelAngle[font=\scriptsize,pos=.8](C,B,A){$\ang{80}$}
		\tkzLabelAngle[font=\scriptsize,pos=.8](A,C,B){$\ang{45}$}
		\tkzLabelSegment[left, font=\scriptsize](A,C){\SI{9}{\centi\meter}} 
		%		\tkzLabelSegment[below, font=\scriptsize](A,B){$7$}
		%		\tkzLabelSegment[above, font=\scriptsize](C,B){$9$}
		\begin{scope}[xshift=4.1cm] 
			\tkzDefPoint(0.5,0){A} 
			\tkzDefShiftPoint[A](+13:2.1){B}
			\tkzDefTriangle[two angles=55 and 80](A,B)\tkzGetPoint{C}  
			\tkzDrawPolygon[line width=1.2pt](A,B,C)
			%		\tkzLabelPoints[font=\scriptsize](A,B,C)
			%		\tkzMarkRightAngle[](B,A,C)
			\tkzMarkAngle[mark=none,-, >=latex, size =0.5](B,A,C)
			\tkzMarkAngle[mark=none,-, >=latex, size =0.5](C,B,A)
			\tkzMarkAngle[mark=none,-, >=latex, size =0.5](A,C,B)
			\tkzLabelAngle[font=\scriptsize,pos=.8](B,A,C){$\ang{55}$}
			\tkzLabelAngle[font=\scriptsize,pos=.8](C,B,A){$\ang{80}$}
			\tkzLabelAngle[font=\scriptsize,pos=.8](A,C,B){$\ang{45}$}
			%		\tkzLabelSegment[left, font=\scriptsize](A,C){$9$} 
			\tkzLabelSegment[below, font=\scriptsize](A,B){\SI{9}{\centi\meter}}
			%		\tkzLabelSegment[above, font=\scriptsize](C,B){$9$}
		\end{scope}
	\end{tikzpicture}   
}
%}
%\fbox{
\parbox[position=t]{0.45\textwidth}{
	%\begin{figure}[H]\centering
	\begin{tikzpicture}[scale=1.0]
		\tkzText(0.5,2.8){$(D)$}
		\tkzSetUpPoint[size=3, fill=black]
		\tkzfctset{tan style/.style={->,>=latex, color=red, line width=1.2pt}}   
		\tkzSetUpLine[line width= 1.2pt, add=.5 and .5] 
		\tkzDefPoint(0.5,0){A} 
		\tkzDefShiftPoint[A](0:4.1){B}
		\tkzDefTriangle[two angles=25 and 125](A,B)\tkzGetPoint{C}  
		\tkzDrawPolygon[line width=1.2pt](A,B,C)
		%		\tkzLabelPoints[font=\scriptsize](A,B,C)
		%		\tkzMarkRightAngle[](B,A,C)
		\tkzMarkAngle[mark=none,-, >=latex, size =0.6](B,A,C)
		\tkzMarkAngle[mark=none,-, >=latex, size =0.4](C,B,A)
		%		\tkzMarkAngle[mark=none,-, >=latex, size =0.5](A,C,B)
		\tkzLabelAngle[font=\scriptsize,pos=1.2](B,A,C){$\ang{20}$}
		\tkzLabelAngle[font=\scriptsize,pos=0.6](C,B,A){$\ang{125}$}
		%		\tkzLabelAngle[font=\scriptsize,pos=.8](A,C,B){$\ang{45}$}
		\tkzLabelSegment[above, font=\scriptsize](A,C){ \SI{13}{\centi\meter}} 
		%		\tkzLabelSegment[below, font=\scriptsize](A,B){$7$}
		%		\tkzLabelSegment[above, font=\scriptsize](C,B){$9$}
		\begin{scope}[xshift=5.0cm] 
			\tkzDefPoint(0.5,0){A} 
			\tkzDefShiftPoint[A](+20:3.1){B}
			\tkzDefTriangle[two angles=25 and 125](A,B)\tkzGetPoint{C}  
			\tkzDrawPolygon[line width=1.2pt](A,B,C)
			%		\tkzLabelPoints[font=\scriptsize](A,B,C)
			%		\tkzMarkRightAngle[](B,A,C)
			\tkzMarkAngle[mark=none,-, >=latex, size =0.7](B,A,C)
			%			\tkzMarkAngle[mark=none,-, >=latex, size =0.5](C,B,A)
			\tkzMarkAngle[mark=none,-, >=latex, size =0.7](A,C,B)
			\tkzLabelAngle[font=\scriptsize,pos=1.3](B,A,C){$\ang{20}$}
			%\tkzLabelAngle[font=\scriptsize,pos=.8](C,B,A){$\ang{35}$}
			\tkzLabelAngle[font=\scriptsize,pos=1.1](A,C,B){$\ang{35}$}
			%		\tkzLabelSegment[left, font=\scriptsize](A,C){$9$} 
			%			\tkzLabelSegment[below, font=\scriptsize](A,B){$9$}
			\tkzLabelSegment[right, font=\scriptsize](C,B){\SI{8}{\centi\meter}}
		\end{scope}
	\end{tikzpicture}   
}
%}


\noindent
%\fbox{
\parbox[position=t]{0.45\textwidth}{
	%\begin{figure}[H]\centering
	\begin{tikzpicture}[scale=1.0]
		\tkzText(0.5,2.8){$(E)$}
		\tkzSetUpPoint[size=3, fill=black]
		\tkzfctset{tan style/.style={->,>=latex, color=red, line width=1.2pt}}   
		\tkzSetUpLine[line width= 1.2pt, add=.5 and .5] 
		\tkzDefPoint(0.5,0){A} 
		\tkzDefShiftPoint[A](0:4.1){B}
		\tkzDefTriangle[two angles=25 and 125](A,B)\tkzGetPoint{C}  
		\tkzDrawPolygon[line width=1.2pt](A,B,C)
		%		\tkzLabelPoints[font=\scriptsize](A,B,C)
		%		\tkzMarkRightAngle[](B,A,C)
		\tkzMarkAngle[mark=none,-, >=latex, size =0.6](B,A,C)
		\tkzMarkAngle[mark=none,-, >=latex, size =0.4](C,B,A)
		%		\tkzMarkAngle[mark=none,-, >=latex, size =0.5](A,C,B)
		%		\tkzLabelAngle[font=\scriptsize,pos=1.2](B,A,C){$\ang{20}$}
		\tkzLabelAngle[font=\scriptsize,pos=0.6](C,B,A){$\ang{135}$}
		%		\tkzLabelAngle[font=\scriptsize,pos=.8](A,C,B){$\ang{45}$}
		%		\tkzLabelSegment[above, font=\scriptsize](A,C){ \SI{13}{\centi\meter}} 
		\tkzLabelSegment[below, font=\scriptsize](A,B){\SI{15}{\centi\meter}}
		\tkzLabelSegment[below right, font=\scriptsize](C,B){\SI{8}{\centi\meter}}
		\begin{scope}[xshift=4.5cm, yshift=-1cm] 
			\tkzDefPoint(0.5,0){A} 
			\tkzDefShiftPoint[A](+20:3.1){B}
			\tkzDefTriangle[two angles=25 and 115](A,B)\tkzGetPoint{C}  
			\tkzDrawPolygon[line width=1.2pt](A,B,C)
			%		\tkzLabelPoints[font=\scriptsize](A,B,C)
			%		\tkzMarkRightAngle[](B,A,C)
			\tkzMarkAngle[mark=none,-, >=latex, size =0.7](B,A,C)
			%			\tkzMarkAngle[mark=none,-, >=latex, size =0.5](C,B,A)
			\tkzMarkAngle[mark=none,-, >=latex, size =0.7](A,C,B)
			\tkzLabelAngle[font=\scriptsize,pos=1.3](B,A,C){$\ang{20}$}
			%\tkzLabelAngle[font=\scriptsize,pos=.8](C,B,A){$\ang{35}$}
			\tkzLabelAngle[font=\scriptsize,pos=1.1](A,C,B){$\ang{25}$}
			%		\tkzLabelSegment[left, font=\scriptsize](A,C){$9$} 
			\tkzLabelSegment[below right, font=\scriptsize](A,B){\SI{15}{\centi\meter}}
			\tkzLabelSegment[right, font=\scriptsize](C,B){\SI{8}{\centi\meter}}
		\end{scope}
	\end{tikzpicture}   
}
%}
\hfill 
%\fbox{
\parbox[position=t]{0.45\textwidth}{
	%\begin{figure}[H]\centering
	\begin{tikzpicture}[scale=1.0]
		\tkzText(0.5,2.5){$(F)$}
		\tkzSetUpPoint[size=3, fill=black]
		\tkzfctset{tan style/.style={->,>=latex, color=red, line width=1.2pt}}   
		\tkzSetUpLine[line width= 1.2pt, add=.5 and .5] 
		\tkzDefPoint(0.5,0){A} 
		\tkzDefShiftPoint[A](-20:3.1){B}
		\tkzDefTriangle[two angles=55 and 70](A,B)\tkzGetPoint{C}  
		\tkzDrawPolygon[line width=1.2pt](A,B,C)
		%		\tkzLabelPoints[font=\scriptsize](A,B,C)
		%		\tkzMarkRightAngle[](B,A,C)
		%		\tkzMarkAngle[mark=none,-, >=latex, size =0.6](B,A,C)
		\tkzMarkAngle[mark=none,-, >=latex, size =0.4](C,B,A)
		\tkzMarkAngle[mark=none,-, >=latex, size =0.5](A,C,B)
		%		\tkzLabelAngle[font=\scriptsize,pos=1.2](B,A,C){$\ang{20}$}
		\tkzLabelAngle[font=\scriptsize,pos=0.6](C,B,A){$\ang{70}$}
		\tkzLabelAngle[font=\scriptsize,pos=.8](A,C,B){$\ang{55}$}
		%		\tkzLabelSegment[above, font=\scriptsize](A,C){ \SI{13}{\centi\meter}} 
		\tkzLabelSegment[below, font=\scriptsize](A,B){\SI{10}{\centi\meter}}
		%\tkzLabelSegment[below right, font=\scriptsize](C,B){\SI{8}{\centi\meter}}
		\begin{scope}[xshift=4.5cm, yshift=-1cm] 
			\tkzDefPoint(0.5,0){A} 
			\tkzDefShiftPoint[A](00:3.5){B}
			\tkzDefTriangle[two angles=55 and 55](A,B)\tkzGetPoint{C}  
			\tkzDrawPolygon[line width=1.2pt](A,B,C)
			%		\tkzLabelPoints[font=\scriptsize](A,B,C)
			%		\tkzMarkRightAngle[](B,A,C)
			\tkzMarkAngle[mark=none,-, >=latex, size =0.6](B,A,C)
			%			\tkzMarkAngle[mark=none,-, >=latex, size =0.5](C,B,A)
			\tkzMarkAngle[mark=none,-, >=latex, size =0.6](A,C,B)
			\tkzLabelAngle[font=\scriptsize,pos=0.9](B,A,C){$\ang{55}$}
			%\tkzLabelAngle[font=\scriptsize,pos=.8](C,B,A){$\ang{35}$}
			\tkzLabelAngle[font=\scriptsize,pos=0.9](A,C,B){$\ang{70}$}
			%		\tkzLabelSegment[left, font=\scriptsize](A,C){$9$} 
			\tkzLabelSegment[below  , font=\scriptsize](A,B){\SI{10}{\centi\meter}}
			%	\tkzLabelSegment[right, font=\scriptsize](C,B){\SI{8}{\centi\meter}}
		\end{scope}
	\end{tikzpicture}   
}
%}


\noindent
\hfill 
%\fbox{
\parbox[position=t]{0.45\textwidth}{
	%\begin{figure}[H]\centering
	\begin{tikzpicture}[scale=1.0]
		\tkzText(0.5,2.5){$(G)$}
		\tkzSetUpPoint[size=3, fill=black]
		\tkzfctset{tan style/.style={->,>=latex, color=red, line width=1.2pt}}   
		\tkzSetUpLine[line width= 1.2pt, add=.5 and .5] 
		\tkzDefPoint(0.5,0){A} 
		\tkzDefShiftPoint[A](-20:2.5){B}
		\tkzDefTriangle[two angles=90 and 55](A,B)\tkzGetPoint{C}  
		\tkzDrawPolygon[line width=1.2pt](A,B,C)
		%		\tkzLabelPoints[font=\scriptsize](A,B,C)
		\tkzMarkRightAngle[](B,A,C)
		%		\tkzMarkAngle[mark=none,-, >=latex, size =0.6](B,A,C)
		%		\tkzMarkAngle[mark=none,-, >=latex, size =0.4](C,B,A)
		%		\tkzMarkAngle[mark=none,-, >=latex, size =0.5](A,C,B)
		%		\tkzLabelAngle[font=\scriptsize,pos=1.2](B,A,C){$\ang{20}$}
		%		\tkzLabelAngle[font=\scriptsize,pos=0.6](C,B,A){$\ang{70}$}
		%		\tkzLabelAngle[font=\scriptsize,pos=.8](A,C,B){$\ang{55}$}
		\tkzLabelSegment[left, font=\scriptsize](A,C){ \SI{9}{\centi\meter}} 
		%	\tkzLabelSegment[below, font=\scriptsize](A,B){\SI{10}{\centi\meter}}
		\tkzLabelSegment[above right, font=\scriptsize](C,B){\SI{12}{\centi\meter}}
		\begin{scope}[xshift=3.5cm, yshift=-1cm] 
			\tkzDefPoint(0.5,0){A} 
			\tkzDefShiftPoint[A](30:3.1){B}
			\tkzDefTriangle[two angles=35 and 90](A,B)\tkzGetPoint{C}  
			\tkzDrawPolygon[line width=1.2pt](A,B,C)
			%		\tkzLabelPoints[font=\scriptsize](A,B,C)
			%					\tkzMarkRightAngle[](B,A,C)
			\tkzMarkRightAngle[](C,B,A) 
			%			\tkzMarkAngle[mark=none,-, >=latex, size =0.6](B,A,C)
			%			\tkzMarkAngle[mark=none,-, >=latex, size =0.5](C,B,A)
			%			\tkzMarkAngle[mark=none,-, >=latex, size =0.6](A,C,B)
			%			\tkzLabelAngle[font=\scriptsize,pos=0.9](B,A,C){$\ang{55}$}
			%\tkzLabelAngle[font=\scriptsize,pos=.8](C,B,A){$\ang{35}$}
			%			\tkzLabelAngle[font=\scriptsize,pos=0.9](A,C,B){$\ang{70}$}
			\tkzLabelSegment[above left, font=\scriptsize](A,C){\SI{12}{\centi\meter}} 
			% \tkzLabelSegment[below  , font=\scriptsize](A,B){\SI{10}{\centi\meter}}
			\tkzLabelSegment[above right, font=\scriptsize](C,B){\SI{8}{\centi\meter}}
		\end{scope}
	\end{tikzpicture}   
}
%}
\hfill 
%\fbox{
\parbox[position=t]{0.45\textwidth}{
	%\begin{figure}[H]\centering
	\begin{tikzpicture}[scale=1.0]
		\tkzText(0.5,2.1){$(H)$}
		\tkzSetUpPoint[size=3, fill=black]
		\tkzfctset{tan style/.style={->,>=latex, color=red, line width=1.2pt}}   
		\tkzSetUpLine[line width= 1.2pt, add=.5 and .5] 
		\tkzDefPoint(0.5,0){A} 
		\tkzDefShiftPoint[A](0:3.8){B}
		\tkzDefTriangle[two angles=-25 and -115](A,B)\tkzGetPoint{C}  
		\tkzDrawPolygon[line width=1.2pt](A,B,C)
		%		\tkzLabelPoints[font=\scriptsize](A,B,C)
		%		\tkzMarkRightAngle[](B,A,C)
		\tkzMarkAngle[mark=none,-, >=latex, size =0.6](C,A,B)
		\tkzMarkAngle[mark=none,-, >=latex, size =0.4](A,B,C)
		%\tkzMarkAngle[mark=none,-, >=latex, size =0.5](A,C,B)
		\tkzLabelAngle[font=\scriptsize,pos=1.2](C,A,B){$\ang{25}$}
		\tkzLabelAngle[font=\scriptsize,pos=0.6](A,B,C){$\ang{115}$}
		%		\tkzLabelAngle[font=\scriptsize,pos=.8](A,C,B){$\ang{55}$}
		%\tkzLabelSegment[above, font=\scriptsize](A,C){ \SI{13}{\centi\meter}} 
		\tkzLabelSegment[above, font=\scriptsize](A,B){\SI{8}{\centi\meter}}
		\tkzLabelSegment[above right, font=\scriptsize, sloped](C,B){\SI{9}{\centi\meter}}
		\begin{scope}[xshift=5.5cm, yshift=-1cm] 
			\tkzDefPoint(0.5,0){A} 
			\tkzDefShiftPoint[A](-10:3.0){B}
			\tkzDefTriangle[two angles=115 and 30](A,B)\tkzGetPoint{C}  
			\tkzDrawPolygon[line width=1.2pt](A,B,C)
			%		\tkzLabelPoints[font=\scriptsize](A,B,C)
			%		\tkzMarkRightAngle[](B,A,C)
			\tkzMarkAngle[mark=none,-, >=latex, size =0.5](B,A,C)
			%			\tkzMarkAngle[mark=none,-, >=latex, size =0.5](C,B,A)
			\tkzMarkAngle[mark=none,-, >=latex, size =0.6](A,C,B)
			\tkzLabelAngle[font=\scriptsize,pos=0.8](B,A,C){$\ang{115}$}
			%	\tkzLabelAngle[font=\scriptsize,pos=.8](C,B,A){$\ang{45}$}
			\tkzLabelAngle[font=\scriptsize,pos=0.9](A,C,B){$\ang{45}$}
			%	\tkzLabelSegment[left, font=\scriptsize](A,C){\SI{9}{\centi\meter}} 
			\tkzLabelSegment[below  , font=\scriptsize](A,B){\SI{8}{\centi\meter}}
			%		\tkzLabelSegment[right, font=\scriptsize](C,B){\SI{8}{\centi\meter}}
		\end{scope}
	\end{tikzpicture}   
}
%}

\noindent
%\fbox{
\parbox[position=t]{0.45\textwidth}{
	%\begin{figure}[H]\centering
	\begin{tikzpicture}[scale=1.0]
		\tkzText(0.5,2.5){$(I)$}
		\tkzSetUpPoint[size=3, fill=black]
		\tkzfctset{tan style/.style={->,>=latex, color=red, line width=1.2pt}}   
		\tkzSetUpLine[line width= 1.2pt, add=.5 and .5] 
		\tkzDefPoint(0.5,0){A} 
		\tkzDefShiftPoint[A](-20:3.1){B}
		\tkzDefTriangle[two angles=55 and 70](A,B)\tkzGetPoint{C}  
		\tkzDrawPolygon[line width=1.2pt](A,B,C)
		%		\tkzLabelPoints[font=\scriptsize](A,B,C)
		%		\tkzMarkRightAngle[](B,A,C)
		\tkzMarkAngle[mark=none,-, >=latex, size =0.5](B,A,C)
		\tkzMarkAngle[mark=none,-, >=latex, size =0.5](C,B,A)
		\tkzMarkAngle[mark=none,-, >=latex, size =0.5](A,C,B)
		\tkzLabelAngle[font=\scriptsize,pos= .8](B,A,C){$\ang{55}$}
		\tkzLabelAngle[font=\scriptsize,pos=0.8](C,B,A){$\ang{80}$}
		\tkzLabelAngle[font=\scriptsize,pos=0.8](A,C,B){$\ang{45}$}
		%		\tkzLabelSegment[above, font=\scriptsize](A,C){ \SI{13}{\centi\meter}} 
		%\tkzLabelSegment[below, font=\scriptsize](A,B){\SI{10}{\centi\meter}}
		\tkzLabelSegment[  above, font=\scriptsize, sloped](C,B){\SI{7}{\centi\meter}}
		\begin{scope}[xshift=4.0cm, yshift=2cm] 
			\tkzDefPoint(0.5,0){A} 
			\tkzDefShiftPoint[A](-5:2.8){B}
			\tkzDefTriangle[two angles=-45 and -80](A,B)\tkzGetPoint{C}  
			\tkzDrawPolygon[line width=1.2pt](A,B,C)
			%		\tkzLabelPoints[font=\scriptsize](A,B,C)
			%		\tkzMarkRightAngle[](B,A,C)
			\tkzMarkAngle[mark=none,-, >=latex, size =0.5](C,A,B)
			\tkzMarkAngle[mark=none,-, >=latex, size =0.5](A,B,C)
			\tkzMarkAngle[mark=none,-, >=latex, size =0.5](B,C,A)
			\tkzLabelAngle[font=\scriptsize,pos= .8](C,A,B){$\ang{45}$}
			\tkzLabelAngle[font=\scriptsize,pos=0.8](A,B,C){$\ang{80}$}
			\tkzLabelAngle[font=\scriptsize,pos=0.8](B,C,A){$\ang{55}$}
			%		\tkzLabelSegment[above, font=\scriptsize](A,C){ \SI{13}{\centi\meter}} 
			\tkzLabelSegment[above, font=\scriptsize](A,B){\SI{7}{\centi\meter}}
			%	\tkzLabelSegment[  above, font=\scriptsize, sloped](C,B){\SI{7}{\centi\meter}}
		\end{scope}
	\end{tikzpicture}   
}
%}
\hfill 
%\fbox{
\parbox[position=t]{0.45\textwidth}{
	%\begin{figure}[H]\centering
	\begin{tikzpicture}[scale=1.0]
		\tkzText(0.5,2.5){$(J)$}
		\tkzSetUpPoint[size=3, fill=black]
		\tkzfctset{tan style/.style={->,>=latex, color=red, line width=1.2pt}}   
		\tkzSetUpLine[line width= 1.2pt, add=.5 and .5] 
		\tkzDefPoint(0.5,0){A} 
		\tkzDefShiftPoint[A](-20:2.5){B}
		\tkzDefTriangle[two angles=80 and 50](A,B)\tkzGetPoint{C}  
		\tkzDrawPolygon[line width=1.2pt](A,B,C)
		%		\tkzLabelPoints[font=\scriptsize](A,B,C)
		%		\tkzMarkRightAngle[](B,A,C)
		%		\tkzMarkAngle[mark=none,-, >=latex, size =0.6](B,A,C)
		\tkzMarkAngle[mark=none,-, >=latex, size =0.4](C,B,A)
		\tkzMarkAngle[mark=none,-, >=latex, size =0.5](A,C,B)
		%		\tkzLabelAngle[font=\scriptsize,pos=1.2](B,A,C){$\ang{20}$}
		\tkzLabelAngle[font=\scriptsize,pos=0.6](C,B,A){$\ang{50}$}
		\tkzLabelAngle[font=\scriptsize,pos=.8](A,C,B){$\ang{40}$}
		\tkzLabelSegment[left, font=\scriptsize](A,C){ \SI{9}{\centi\meter}} 
		%	\tkzLabelSegment[below, font=\scriptsize](A,B){\SI{10}{\centi\meter}}
		%		\tkzLabelSegment[above right, font=\scriptsize](C,B){\SI{12}{\centi\meter}}
		\begin{scope}[xshift=3.5cm, yshift=+2cm] 
			\tkzDefPoint(0.5,0){A} 
			\tkzDefShiftPoint[A](-60:3.1){B}
			\tkzDefTriangle[two angles=90 and 35](A,B)\tkzGetPoint{C}  
			\tkzDrawPolygon[line width=1.2pt](A,B,C)
			%		\tkzLabelPoints[font=\scriptsize](A,B,C)
			\tkzMarkRightAngle[](B,A,C)
			%			\tkzMarkRightAngle[](C,B,A) 
			%			\tkzMarkAngle[mark=none,-, >=latex, size =0.6](B,A,C)
			%			\tkzMarkAngle[mark=none,-, >=latex, size =0.5](C,B,A)
			%			\tkzMarkAngle[mark=none,-, >=latex, size =0.6](A,C,B)
			%			\tkzLabelAngle[font=\scriptsize,pos=0.9](B,A,C){$\ang{55}$}
			%\tkzLabelAngle[font=\scriptsize,pos=.8](C,B,A){$\ang{35}$}
			%			\tkzLabelAngle[font=\scriptsize,pos=0.9](A,C,B){$\ang{70}$}
			\tkzLabelSegment[above left, font=\scriptsize](A,C){\SI{9}{\centi\meter}} 
			% \tkzLabelSegment[below  , font=\scriptsize](A,B){\SI{10}{\centi\meter}}
			\tkzLabelSegment[above right, font=\scriptsize](C,B){\SI{12}{\centi\meter}}
		\end{scope}
	\end{tikzpicture}   
}
%}



\noindent
%\fbox{
\parbox[position=t]{0.45\textwidth}{
	%\begin{figure}[H]\centering
	\begin{tikzpicture}[scale=1.0]
		\tkzText(0.5,2.8){$(K)$}
		\tkzSetUpPoint[size=3, fill=black]
		\tkzfctset{tan style/.style={->,>=latex, color=red, line width=1.2pt}}   
		\tkzSetUpLine[line width= 1.2pt, add=.5 and .5] 
		\tkzDefPoint(0.5,0){A} 
		\tkzDefShiftPoint[A](-15:2.8){B}
		\tkzDefTriangle[two angles=55 and 80](A,B)\tkzGetPoint{C}  
		\tkzDrawPolygon[line width=1.2pt](A,B,C)
		%		\tkzLabelPoints[font=\scriptsize](A,B,C)
		%		\tkzMarkRightAngle[](B,A,C)
		\tkzMarkAngle[mark=none,-, >=latex, size =0.5](B,A,C)
		\tkzMarkAngle[mark=none,-, >=latex, size =0.5](C,B,A)
		\tkzMarkAngle[mark=none,-, >=latex, size =0.5](A,C,B)
		\tkzLabelAngle[font=\scriptsize,pos=.8](B,A,C){$\ang{55}$}
		\tkzLabelAngle[font=\scriptsize,pos=.8](C,B,A){$\ang{80}$}
		\tkzLabelAngle[font=\scriptsize,pos=.8](A,C,B){$\ang{45}$}
		\tkzLabelSegment[left, font=\scriptsize](A,C){\SI{4}{\centi\meter}} 
		\tkzLabelSegment[below, font=\scriptsize](A,B){\SI{2.5}{\centi\meter}}
		%		\tkzLabelSegment[above, font=\scriptsize](C,B){$9$}
		\begin{scope}[xshift=4.1cm] 
			\tkzDefPoint(0.5,0){A} 
			\tkzDefShiftPoint[A](+13:2.3){B}
			\tkzDefTriangle[two angles=80 and 55](A,B)\tkzGetPoint{C}  
			\tkzDrawPolygon[line width=1.2pt](A,B,C)
			%		\tkzLabelPoints[font=\scriptsize](A,B,C)
			%		\tkzMarkRightAngle[](B,A,C)
			\tkzMarkAngle[mark=none,-, >=latex, size =0.5](B,A,C)
			\tkzMarkAngle[mark=none,-, >=latex, size =0.5](C,B,A)
			\tkzMarkAngle[mark=none,-, >=latex, size =0.5](A,C,B)
			\tkzLabelAngle[font=\scriptsize,pos=.8](B,A,C){$\ang{80}$}
			\tkzLabelAngle[font=\scriptsize,pos=.8](C,B,A){$\ang{55}$}
			\tkzLabelAngle[font=\scriptsize,pos=.8](A,C,B){$\ang{45}$}
			%		\tkzLabelSegment[left, font=\scriptsize](A,C){$9$} 
			\tkzLabelSegment[below, font=\scriptsize](A,B){\SI{2.5}{\centi\meter}}
			\tkzLabelSegment[right, font=\scriptsize](C,B){\SI{4}{\centi\meter}}
		\end{scope}
	\end{tikzpicture}   
}
%}
\hfill 
%\fbox{
\parbox[position=t]{0.45\textwidth}{
	%\begin{figure}[H]\centering
	\begin{tikzpicture}[scale=1.0]
		\tkzText(0.5,2.8){$(L)$}
		\tkzSetUpPoint[size=3, fill=black]
		\tkzfctset{tan style/.style={->,>=latex, color=red, line width=1.2pt}}   
		\tkzSetUpLine[line width= 1.2pt, add=.5 and .5] 
		\tkzDefPoint(1,0){A} 
		\tkzDefShiftPoint[A](0:3.3){B}
		\tkzDefTriangle[two angles=90 and 35](A,B)\tkzGetPoint{C}  
		\tkzDrawPolygon[line width=1.2pt](A,B,C)
		%		\tkzLabelPoints[font=\scriptsize](A,B,C)
		\tkzMarkRightAngle[](B,A,C)
		%\tkzMarkAngle[mark=none,-, >=latex, size =0.5](C,B,A)
		\tkzMarkAngle[mark=none,-, >=latex, size =0.5](A,C,B)
		%\tkzLabelAngle[font=\scriptsize,pos=.8](C,B,A){$\ang{70}$}
		\tkzLabelAngle[font=\scriptsize,pos=.8](A,C,B){$\ang{70}$}
		%\tkzLabelSegment[above, font=\scriptsize](A,C){$10$} 
		\tkzLabelSegment[below, font=\scriptsize](A,B){\SI{8}{\centi\meter}}
		%		\tkzLabelSegment[above, font=\scriptsize](C,B){\SI{9}{\centi\meter}}
		\begin{scope}[xshift=5cm] 
			\tkzDefPoint(0,0){A}
			\tkzDefShiftPoint[A](+45:3.5){B} 
			\tkzDefTriangle[two angles=25 and 90](A,B)\tkzGetPoint{C}  
			
			\tkzDrawPolygon[line width=1.2pt](A,B,C)
			%	\tkzLabelPoints[font=\scriptsize](A,B,C)
			%			\tkzMarkRightAngle[](B,A,C)
			\tkzMarkRightAngle[](C,B,A)
			%\tkzMarkAngle[mark=none,-, >=latex, size =0.5](C,B,A)
			\tkzMarkAngle[mark=none,-, >=latex, size =0.5](A,C,B)
			%\tkzLabelAngle[font=\scriptsize,pos=.8](C,B,A){$\ang{70}$}
			\tkzLabelAngle[font=\scriptsize,pos=.8](A,C,B){$\ang{70}$}
			\tkzLabelSegment[below right, font=\scriptsize](A,B){\SI{8}{\centi\meter}} 
			%	\tkzLabelSegment[left, font=\scriptsize](A,C){\SI{7}{\centi\meter}}
			% 	\tkzLabelSegment[above right, font=\scriptsize](C,B){\SI{9}{\centi\meter}}
		\end{scope}
	\end{tikzpicture}   
}
%}

\noindent
%\fbox{
\parbox[position=t]{0.45\textwidth}{
	%\begin{figure}[H]\centering
	\begin{tikzpicture}[scale=1.0]
		\tkzText(0.5,2.1){$(M)$}
		\tkzSetUpPoint[size=3, fill=black]
		\tkzfctset{tan style/.style={->,>=latex, color=red, line width=1.2pt}}   
		\tkzSetUpLine[line width= 1.2pt, add=.5 and .5] 
		\tkzDefPoint(0.5,0){A} 
		\tkzDefShiftPoint[A](0:3.5){B}
		\tkzDefTriangle[two angles=-20 and -130](A,B)\tkzGetPoint{C}  
		\tkzDrawPolygon[line width=1.2pt](A,B,C)
		%		\tkzLabelPoints[font=\scriptsize](A,B,C)
		%		\tkzMarkRightAngle[](B,A,C)
		\tkzMarkAngle[mark=none,-, >=latex, size =0.8](C,A,B)
		\tkzMarkAngle[mark=none,-, >=latex, size =0.4](A,B,C)
		\tkzMarkAngle[mark=none,-, >=latex, size =0.6](B,C,A)
		\tkzLabelAngle[font=\scriptsize,pos=1.3](C,A,B){$\ang{20}$}
		\tkzLabelAngle[font=\scriptsize,pos=0.6](A,B,C){$\ang{130}$}
		\tkzLabelAngle[font=\scriptsize,pos=1.2](B,C,A){$\ang{30}$}
		%\tkzLabelSegment[above, font=\scriptsize](A,C){ \SI{13}{\centi\meter}} 
		%		\tkzLabelSegment[above, font=\scriptsize](A,B){\SI{8}{\centi\meter}}
		%		\tkzLabelSegment[above right, font=\scriptsize, sloped](C,B){\SI{9}{\centi\meter}}
		\begin{scope}[xshift=4.5cm, yshift= -0.5cm] 
			\tkzDefPoint(0.5,0){A} 
			\tkzDefShiftPoint[A](-5:3.5){B}
			\tkzDefTriangle[two angles=130 and 20](A,B)\tkzGetPoint{C}  
			\tkzDrawPolygon[line width=1.2pt](A,B,C)
			%		\tkzLabelPoints[font=\scriptsize](A,B,C)
			%		\tkzMarkRightAngle[](B,A,C)
			\tkzMarkAngle[mark=none,-, >=latex, size =0.4](B,A,C)
			\tkzMarkAngle[mark=none,-, >=latex, size =0.5](C,B,A)
			\tkzMarkAngle[mark=none,-, >=latex, size =0.8](A,C,B)
			\tkzLabelAngle[font=\scriptsize,pos=0.8](B,A,C){$\ang{130}$}
			\tkzLabelAngle[font=\scriptsize,pos=1.3](C,B,A){$\ang{20}$}
			\tkzLabelAngle[font=\scriptsize,pos=1.2](A,C,B){$\ang{30}$}
			%	\tkzLabelSegment[left, font=\scriptsize](A,C){\SI{9}{\centi\meter}} 
			%			\tkzLabelSegment[below  , font=\scriptsize](A,B){\SI{8}{\centi\meter}}
			%		\tkzLabelSegment[right, font=\scriptsize](C,B){\SI{8}{\centi\meter}}
		\end{scope}
	\end{tikzpicture}   
}
%}
\hfill 
%\fbox{
\parbox[position=t]{0.45\textwidth}{
	%\begin{figure}[H]\centering
	\begin{tikzpicture}[scale=1.0]
		\tkzText(0.5,2.1){$(N)$}
		\tkzSetUpPoint[size=3, fill=black]
		\tkzfctset{tan style/.style={->,>=latex, color=red, line width=1.2pt}}   
		\tkzSetUpLine[line width= 1.2pt, add=.5 and .5] 
		\tkzDefPoint(0.5,0){A} 
		\tkzDefShiftPoint[A](0:2.5){B}
		\tkzDefTriangle[two angles=-45 and -90](A,B)\tkzGetPoint{C}  
		\tkzDrawPolygon[line width=1.2pt](A,B,C)
		%		\tkzLabelPoints[font=\scriptsize](A,B,C)
		%			\tkzMarkRightAngle[](B,A,C)
		\tkzMarkRightAngle[](A,B,C)
		%	\tkzMarkAngle[mark=none,-, >=latex, size =0.6](C,A,B)
		%	\tkzMarkAngle[mark=none,-, >=latex, size =0.4](A,B,C)
		%\tkzMarkAngle[mark=none,-, >=latex, size =0.5](A,C,B)
		%	\tkzLabelAngle[font=\scriptsize,pos=1.2](C,A,B){$\ang{25}$}
		%	\tkzLabelAngle[font=\scriptsize,pos=0.6](A,B,C){$\ang{115}$}
		%		\tkzLabelAngle[font=\scriptsize,pos=.8](A,C,B){$\ang{55}$}
		%\tkzLabelSegment[above, font=\scriptsize](A,C){ \SI{13}{\centi\meter}} 
		%	\tkzLabelSegment[above, font=\scriptsize](A,B){\SI{8}{\centi\meter}}
		%	\tkzLabelSegment[above right, font=\scriptsize, sloped](C,B){\SI{9}{\centi\meter}}
		\tkzMarkSegments[mark=s||](A,B B,C)
		\begin{scope}[xshift=4.0cm, yshift=-1cm] 
			\tkzDefPoint(0.5,0){A} 
			\tkzDefShiftPoint[A](-10:3.0){B}
			\tkzDefTriangle[two angles=90 and 45](A,B)\tkzGetPoint{C}  
			\tkzDrawPolygon[line width=1.2pt](A,B,C)
			%		\tkzLabelPoints[font=\scriptsize](A,B,C)
			\tkzMarkRightAngle[](B,A,C)
			%		\tkzMarkAngle[mark=none,-, >=latex, size =0.5](B,A,C)
			%			\tkzMarkAngle[mark=none,-, >=latex, size =0.5](C,B,A)
			%		\tkzMarkAngle[mark=none,-, >=latex, size =0.6](A,C,B)
			%		\tkzLabelAngle[font=\scriptsize,pos=0.8](B,A,C){$\ang{115}$}
			%	\tkzLabelAngle[font=\scriptsize,pos=.8](C,B,A){$\ang{45}$}
			%		\tkzLabelAngle[font=\scriptsize,pos=0.9](A,C,B){$\ang{45}$}
			%	\tkzLabelSegment[left, font=\scriptsize](A,C){\SI{9}{\centi\meter}} 
			%		\tkzLabelSegment[below  , font=\scriptsize](A,B){\SI{8}{\centi\meter}}
			%		\tkzLabelSegment[right, font=\scriptsize](C,B){\SI{8}{\centi\meter}}
			\tkzMarkSegments[mark=s||](A,B A,C)
		\end{scope}
	\end{tikzpicture}   
}
%} 

\begin{exercise}\hfill 
	
	$ABCD$ est un parallélogramme. $O$ est l'intersection des diagonales. 
	\begin{enumerate}[label={\alph*)},leftmargin=*]
		\item Montrer que les triangles $ABC$ et $ACD$ sont égaux.
		\item Montrer que les triangles $OAB$ et $OCD$ sont égaux.
	\end{enumerate} 
\end{exercise}

\begin{problem}\hfill 
	
	$B$ est un point du segment $[AC]$.  Les triangles $ABD$ et $BCE$ sont situés du même côté du segment $[AC]$ et sont équilatéraux.  Le point $I$ est l'intersection des segments $[AE]$ et $[CD]$, et $J$ est l'intersection de $[BD]$ et $[AE]$.
	
	\noindent\parbox[position=t]{0.55 \linewidth}{  
		\begin{enumerate}[label={\alph*)},leftmargin=*]
			\item Montrer que les triangles $ABE$ et $BCD$ sont égaux.
			\item En déduire que $AE=CD$ et $\widehat{JAB}=\widehat{JDI}$.
			\item Montrer que  $\widehat{DIJ}=\ang{60}$. 
		\end{enumerate}
	}
	\hfill 
	\parbox[position=t]{0.45 \linewidth}{  
		%	\begin{figure}[H]
		\begin{tikzpicture}[scale=0.7]
			\tkzSetUpPoint[size=2, fill=black]
			\tkzfctset{tan style/.style={->,>=latex, color=red, line width=1.2pt}}   
			\tkzSetUpLine[line width= 1.2pt, add=.5 and .5]  
			%\tkzDefPoints{0/0/A,3/0/B,7/0/C}
			\tkzDefPoint(0,0){A}
			\tkzDefPoint(3,0){B}
			\tkzDefPoint(7,0){C}
			\tkzDefEquilateral(A,B) \tkzGetPoint{D}
			\tkzDefEquilateral(B,C) \tkzGetPoint{E}
			\tkzDrawPolygon[line width = 1pt](A,B,D)
			\tkzDrawPolygon[line width = 1pt](B,C,E)
			\tkzDrawSegments[dashed](D,C  E,A)
			\tkzLabelPoints[below, font=\scriptsize](A,B,C) 
			\tkzInterLL(A,E)(C,D) \tkzGetPoint{I}
			\tkzInterLL(A,E)(B,D) \tkzGetPoint{J}
			\tkzDrawPoints(A,B,C,D,E,I,J) 
			\tkzLabelPoints[above,font=\scriptsize](D,E,I)
			\tkzLabelPoint[left,font=\scriptsize](J){$J$} 
		\end{tikzpicture}
		%	\end{figure}
	}  
\end{problem}
\bigskip 
\begin{problem}\hfill 
	
	
	
	\noindent\parbox[position=t]{0.50 \linewidth}{  
		$B$ est un point du segment $[AC]$. 
		
		$ABDE$ et $CDGF$ sont des carrés. 
		
		$I$ est l'intersection de $(CE)$ et de $(BG)$.
		
		$J$ est l'intersection de $(CE)$ et de $(BD)$.
		\begin{enumerate}[label={\alph*)},leftmargin=*] 
			\item Montrer que $EC=BG$
			\item Montrer que les triangles $EDJ$ et $IJB$ sont semblables.
			\item En déduire que les droites $(CE)$ et $(BG)$ sont perpendiculaires.
		\end{enumerate}
	}
	\hfill 
	\parbox[position=t]{0.45  \linewidth}{  
		%	\begin{figure}[H]
		\begin{tikzpicture}[scale=1.0]
			\tkzSetUpPoint[size=2, fill=black]
			\tkzfctset{tan style/.style={->,>=latex, color=red, line width=1.2pt}}   
			\tkzSetUpLine[line width= 1.2pt, add=.5 and .5]  
			%\tkzDefPoints{0/0/A,3/0/B,7/0/C}
			\tkzDefPoint(0,0){A}
			\tkzDefPoint(2.2,0){B}
			\tkzDefPoint(5.3,0){C}
			\tkzDefSquare(A,B) \tkzGetPoints{D}{E}
			\tkzDefSquare(D,C) \tkzGetPoints{F}{G} 
			\tkzDrawPolygon[line width = 1.2pt](A,B,D,E)
			\tkzDrawPolygon[line width = 1.2pt](D,C,F,G)
			\tkzDrawSegment[line width = 1.2pt](B,C)
			\tkzMarkRightAngle(C,B,D)
			\tkzDrawSegments[dashed](E,C  B,G)
			\tkzLabelPoints[below, font=\scriptsize](A,B,C) 
			\tkzInterLL(E,C)(B,G) \tkzGetPoint{I}
			\tkzInterLL(E,C)(B,D) \tkzGetPoint{J}
			\tkzDrawPoints(A,B,C,D,E,F,G) 
			\tkzLabelPoints[above,font=\scriptsize](E,F,G)
			\tkzLabelPoints[above left,font=\scriptsize](D)
			\tkzLabelPoints[below right ,font=\scriptsize](I)
			\tkzLabelPoint[below left,font=\scriptsize](J){$J$} 
		\end{tikzpicture}
		%	\end{figure}
	}  
\end{problem}
%----------------------------------------------------------------------------------------
%	 
%----------------------------------------------------------------------------------------
\cleardoublepage
\end{document}